\chapter{Matriz de Rastreabilidade Requisitos vs. Modelo de Dados}
\label{cap_matriz_rastreabilidade}

A rastreabilidade bidirecional é um critério essencial de engenharia de software para assegurar que cada requisito funcional especificado possua sustentação formal na camada de persistência e que nenhuma estrutura de banco de dados exista sem justificativa de negócio clara \cite{elmasri2016}.

As Tabelas \ref{tab_rastreabilidade_dados_1} e \ref{tab_rastreabilidade_dados_2} mapeiam exaustivamente os 25 Requisitos Funcionais do SIGEA-GTT em relação às tabelas físicas, colunas centrais e mecanismos de integridade do SGBD:

\begin{table}[!ht]
\centering
\caption{Matriz de Rastreabilidade: Requisitos Funcionais (RF01 a RF13) vs. Modelo de Dados}
\label{tab_rastreabilidade_dados_1}
\footnotesize
\begin{tabularx}{\textwidth}{p{1.2cm} p{3.8cm} p{3.8cm} X}
\toprule
\textbf{Código} & \textbf{Nome do Requisito Funcional} & \textbf{Tabelas Físicas} & \textbf{Atributos e Mecanismos de Integridade} \\
\midrule
RF01 & Autenticação e Domínio Exclusivo & \texttt{users} & \texttt{email}, \texttt{password\_hash}, \texttt{chk\_email\_domain} \\
RF02 & Primeiro Acesso e Troca de Senha & \texttt{users} & \texttt{must\_change\_password}, \texttt{password\_hash}, \texttt{updated\_at} \\
RF03 & Alteração de Senha no Perfil & \texttt{users} & \texttt{password\_hash}, \texttt{updated\_at} \\
RF04 & Perfis de Usuário Estritos & \texttt{users} & \texttt{role}, tipo enumerado \texttt{user\_role} \\
RF05 & Matrícula Exclusiva para Discentes & \texttt{users} & \texttt{registration\_number}, \texttt{chk\_registration\_required} \\
RF06 & Gestão de Módulos GTT & \texttt{gtt\_modules} & \texttt{code}, \texttt{name}, \texttt{description}, \texttt{is\_active} \\
RF07 & Gestão de Gatilhos Clínicos & \texttt{gtt\_triggers} & \texttt{module\_id}, \texttt{code}, \texttt{name}, \texttt{description} \\
RF08 & Catálogo Interativo Didático & \texttt{gtt\_modules}, \texttt{gtt\_triggers} & Índices B-Tree \texttt{idx\_gtt\_modules\_code}, \texttt{idx\_gtt\_triggers\_code} \\
RF09 & Escala de Gravidades NCC MERP & \texttt{harm\_severities} & \texttt{category\_letter}, \texttt{is\_harm}, \texttt{idx\_harm\_severities\_letter} \\
RF10 & Formatação Universal de Turmas & \texttt{academic\_classes} & \texttt{subject\_name}, \texttt{class\_code}, \texttt{academic\_period}, \texttt{uk\_class\_identifier} \\
RF11 & Professor Titular Obrigatório & \texttt{academic\_classes} & \texttt{professor\_id} com \texttt{NOT NULL} e \texttt{RESTRICT} \\
RF12 & Matrícula de Alunos em Turmas & \texttt{class\_students} & \texttt{class\_id}, \texttt{student\_id}, \texttt{enrolled\_at}, PK Composta \\
RF13 & Bloqueio de Fechamento de Turma & \texttt{academic\_classes} & \texttt{is\_closed}, verificação transacional de pendências \\
\bottomrule
\end{tabularx}
\fonte{Elaborado pelo autor (2026).}
\end{table}

\begin{table}[!ht]
\centering
\caption{Matriz de Rastreabilidade: Requisitos Funcionais (RF14 a RF25) vs. Modelo de Dados}
\label{tab_rastreabilidade_dados_2}
\footnotesize
\begin{tabularx}{\textwidth}{p{1.2cm} p{3.8cm} p{3.8cm} X}
\toprule
\textbf{Código} & \textbf{Nome do Requisito Funcional} & \textbf{Tabelas Físicas} & \textbf{Atributos e Mecanismos de Integridade} \\
\midrule
RF14 & Gestão de Atividades de Auditoria & \texttt{activities} & \texttt{class\_id}, \texttt{title}, \texttt{description}, \texttt{deadline} \\
RF15 & Prontuário Simulado Estruturado & \texttt{activities} & Coluna \texttt{clinical\_case\_data} em formato \texttt{JSONB} \\
RF16 & Auditoria sob Cronômetro de 20 min & \texttt{activity\_submissions} & \texttt{submission\_date}, controle de sessão via JWT \\
RF17 & Rastreamento de Gatilhos e Danos & \texttt{activity\_submissions} & Coluna \texttt{identified\_triggers} em formato \texttt{JSONB} \\
RF18 & Diagrama de Ishikawa (6M) & \texttt{activity\_submissions} & Seção \texttt{'ishikawa'} em \texttt{quality\_tools\_data} \\
RF19 & Plano de Ação 5W2H & \texttt{activity\_submissions} & Seção \texttt{'plan\_5w2h'} em \texttt{quality\_tools\_data} \\
RF20 & Matriz de Priorização GUT & \texttt{activity\_submissions} & Seção \texttt{'matrix\_gut'} em \texttt{quality\_tools\_data} \\
RF21 & Ciclo PDCA da Qualidade & \texttt{activity\_submissions} & Seção \texttt{'cycle\_pdca'} em \texttt{quality\_tools\_data} \\
RF22 & Correção Docente e Feedback & \texttt{activity\_submissions} & \texttt{grade}, \texttt{professor\_feedback}, \texttt{graded\_at}, \texttt{chk\_grade\_range} \\
RF23 & Comparação com Gabarito Oficial & \texttt{activity\_submissions} & Cruzamento de \texttt{identified\_triggers} com gabarito da atividade \\
RF24 & Consolidação de Indicadores IHI & Múltiplas tabelas & Agregações SQL e fórmulas IHI-GTT (EAs por 1.000 pacientes-dia) \\
RF25 & Painéis e Relatórios Sintéticos & Múltiplas tabelas & Consultas agregadas otimizadas por índices B-Tree e GIN \\
\bottomrule
\end{tabularx}
\fonte{Elaborado pelo autor (2026).}
\end{table}

A análise das matrizes comprova que 100\% dos requisitos funcionais encontram respaldo formal nas tabelas, colunas, tipos e restrições de integridade do banco de dados, inexistindo lacunas estruturais de persistência.
